% =============================================================================
% Capítulo 4 — Implementação
% =============================================================================
\chapter{Implementação}
\label{cap:implementacao}

Este capítulo descreve os detalhes técnicos da implementação do
\textit{vertical slice}, apresentando a organização do código em módulos, os
sistemas de \textit{gameplay} que compõem o \textit{baseline} e o personagem
jogável, o núcleo de proveniência responsável pelo registro causal das
sessões, a integração com aprendizado por reforço via \textit{Reward Shaping}
e \texttt{BossAgent}, e os sistemas de coleta de métricas e gerenciamento de
sessões.

% -----------------------------------------------------------------------------
\section{Arquitetura do Código}
\label{sec:arquitetura-codigo}

A base de código foi organizada em módulos com responsabilidades claramente
delimitadas, seguindo o princípio de separação de responsabilidades
(\textit{Separation of Concerns}). Cada módulo encapsula uma área funcional
do sistema e expõe sua interface pública por meio de eventos C\#
(\texttt{Action<T>} e \textit{delegates}), minimizando o acoplamento direto
entre componentes.

A \autoref{tab:modulos} apresenta os módulos e suas responsabilidades.

\begin{table}[htb]
    \centering
    \caption{Módulos da arquitetura e suas responsabilidades.}
    \label{tab:modulos}
    \begin{tabular}{l p{9.5cm}}
        \toprule
        \textbf{Módulo}     & \textbf{Responsabilidade} \\
        \midrule
        Player              & Captura de entrada, movimento, ataque, \textit{dash}
                              e sistema de vida do personagem Elian. \\
        \addlinespace
        Boss FSM             & \textit{Baseline} determinístico com estados
                              \textit{Idle}, \textit{Chase}, \textit{Attack} e
                              \textit{Retreat}. \\
        \addlinespace
        Proveniência         & Registro causal de eventos em memória, exportação
                              para JSON e módulo de \textit{Reward Shaping}. \\
        \addlinespace
        Boss PPO             & Agente ML-Agents com observações vetoriais, ações
                              discretas e integração com o
                              \texttt{ProvenanceRewardShaper}. \\
        \addlinespace
        Arena / Métricas     & Ciclo de vida da sessão (início, vitória, derrota),
                              coleta de métricas e exportação CSV. \\
        \bottomrule
    \end{tabular}
    \fonte{Elaborado pelo autor.}
\end{table}

\subsection{Comunicação por Eventos}
\label{subsec:comunicacao-eventos}

A comunicação entre módulos é inteiramente dirigida por eventos. Cada
componente que produz informação relevante --- como a execução de um ataque,
a recepção de dano ou a morte de um personagem --- expõe um evento público
(\texttt{Action} ou \texttt{Action<T>}) ao qual outros componentes podem se
inscrever. Esse padrão arquitetural garante que:

\begin{itemize}
    \item O emissor do evento desconhece seus consumidores, eliminando
          dependências circulares;
    \item Novos sistemas (como o de proveniência ou o de métricas) podem ser
          adicionados sem modificar os scripts de \textit{gameplay} existentes,
          bastando inscrever-se nos eventos já expostos;
    \item A testabilidade dos componentes é preservada, pois cada módulo pode
          ser verificado isoladamente por meio de testes unitários que disparam
          ou escutam eventos sem depender da cena completa do Unity.
\end{itemize}

Por exemplo, o evento \texttt{OnBossAttackPerformed}, emitido pelo
\texttt{BossFsmController} quando o chefe executa um ataque, é consumido
simultaneamente pelo \texttt{ProvenanceLogger} (para registrar o evento no
grafo), pelo \texttt{GameMetrics} (para incrementar a contagem de ataques) e
pelo \texttt{BossAttackDamage} (para aplicar dano ao jogador). Nenhum desses
consumidores conhece os demais.

% -----------------------------------------------------------------------------
\section{Sistemas de Gameplay e Baseline}
\label{sec:sistemas-gameplay}

Esta seção descreve a implementação dos scripts que compõem o personagem
jogável e o chefe \textit{baseline} controlado por FSM.

\subsection{Personagem Jogável --- PlayerController}
\label{subsec:player-controller}

O \texttt{PlayerController} é o componente central do personagem Elian. Ele
captura as entradas do jogador por meio do \textit{Unity Input System} e
traduz essas entradas em ações físicas executadas via \texttt{Rigidbody2D}.
A separação entre captura de \textit{input} e execução de física segue as
diretrizes de arquitetura do projeto: as entradas são lidas no \texttt{Update}
e as forças são aplicadas no \texttt{FixedUpdate}, garantindo consistência
com o \textit{timestep} fixo da simulação física.

As ações implementadas são:

\begin{itemize}
    \item \textbf{Movimento horizontal}: aplica velocidade ao
          \texttt{Rigidbody2D} na direção indicada pelo eixo de entrada;
    \item \textbf{Pulo}: aplica impulso vertical quando o personagem está em
          contato com o solo, verificado por um \textit{ground check}
          via \textit{raycast} ou \textit{overlap};
    \item \textbf{Dash}: desloca o personagem horizontalmente com velocidade
          aumentada por um curto período;
    \item \textbf{Ataque}: delega a lógica de dano ao componente
          \texttt{PlayerCombat}.
\end{itemize}

Cada ação emite um evento público quando executada --- por exemplo,
\texttt{OnPlayerJumped}, \texttt{OnPlayerAttacked} e
\texttt{OnPlayerDashed} --- permitindo que o \texttt{ProvenanceLogger}
registre o evento sem que o \texttt{PlayerController} tenha conhecimento do
sistema de proveniência.

\subsection{Combate --- PlayerCombat e IDamageable}
\label{subsec:player-combat}

O componente \texttt{PlayerCombat} implementa a lógica de detecção e
aplicação de dano do ataque do jogador. A cada execução de ataque, o
componente utiliza o método \texttt{Physics2D.OverlapBoxAll} para detectar
todos os \textit{colliders} dentro de uma área retangular posicionada à
frente do personagem. Os objetos detectados que implementam a interface
\texttt{IDamageable} recebem dano por meio da chamada ao método
\texttt{TakeDamage(float amount)}.

A adoção da interface \texttt{IDamageable} como contrato de dano é uma
decisão de projeto que promove extensibilidade: qualquer entidade do jogo que
precise receber dano --- sejam chefes, inimigos comuns ou objetos destrutíveis
--- pode implementar essa interface sem alterar o código de ataque do jogador.
No contexto do \textit{vertical slice}, tanto o \texttt{BossHealth} (chefe
FSM) quanto o chefe PPO implementam \texttt{IDamageable}, garantindo que o
mesmo sistema de combate funcione em ambas as condições experimentais.

\subsection{Vida do Jogador --- PlayerHealth}
\label{subsec:player-health}

O componente \texttt{PlayerHealth} gerencia os pontos de vida do personagem
Elian. Ele implementa a interface \texttt{IDamageable} e mantém um valor de
vida corrente com máximo configurável (\texttt{maxHealth} = 100). Quando o
valor de vida atinge zero, o componente emite um evento de morte que é
consumido pelo \texttt{ArenaManager} para encerrar a sessão como derrota.

Esse componente foi adicionado durante a preparação para a coleta oficial,
após a identificação de que a ausência de vida real no jogador comprometeriam
a validade das métricas: sem possibilidade de derrota, o campo
\texttt{player\_damage\_taken} seria artificial e o resultado da sessão
tenderia invariavelmente a vitória.

\subsection{Boss FSM --- BossFsmController}
\label{subsec:boss-fsm-impl}

O \texttt{BossFsmController} implementa o chefe \textit{baseline}
como uma máquina de estados finitos com quatro estados, operando sobre um
\texttt{Rigidbody2D}. As transições são determinadas exclusivamente pela
distância horizontal até o jogador e pelo estado interno de
\textit{cooldown}.

A \autoref{tab:fsm-estados} descreve o comportamento associado a cada estado.

\begin{table}[htb]
    \centering
    \caption{Estados e comportamentos do chefe FSM.}
    \label{tab:fsm-estados}
    \begin{tabular}{l p{10cm}}
        \toprule
        \textbf{Estado}       & \textbf{Comportamento} \\
        \midrule
        \textit{Idle}         & O chefe permanece parado. Ocorre tipicamente
                                quando o jogador está dentro do alcance de
                                ataque, mas o \textit{cooldown} ainda não
                                expirou. \\
        \addlinespace
        \textit{Chase}        & O chefe move-se horizontalmente em direção ao
                                jogador com velocidade de perseguição
                                configurável. \\
        \addlinespace
        \textit{Attack}       & O chefe interrompe o movimento, emite o evento
                                \texttt{OnBossAttackPerformed}, exibe
                                \textit{feedback} visual e entra em
                                \textit{cooldown}. \\
        \addlinespace
        \textit{Retreat}      & O chefe move-se na direção oposta ao jogador
                                por um curto período, criando espaçamento
                                após o ataque. \\
        \bottomrule
    \end{tabular}
    \fonte{Elaborado pelo autor.}
\end{table}

Os parâmetros do FSM foram calibrados para produzir um oponente funcional mas
previsível, conforme apresentado na \autoref{tab:fsm-params}.

\begin{table}[htb]
    \centering
    \caption{Parâmetros do chefe FSM.}
    \label{tab:fsm-params}
    \begin{tabular}{l r l}
        \toprule
        \textbf{Parâmetro}          & \textbf{Valor} & \textbf{Unidade} \\
        \midrule
        \texttt{chaseSpeed}         & 2,0     & unidades/s         \\
        \texttt{attackRange}        & 1,25    & unidades           \\
        \texttt{attackDuration}     & 0,2     & segundos           \\
        \texttt{retreatDuration}    & 0,45    & segundos           \\
        \texttt{retreatSpeed}       & 1,5     & unidades/s         \\
        \texttt{attackCooldown}     & 0,8     & segundos           \\
        \bottomrule
    \end{tabular}
    \fonte{Elaborado pelo autor.}
\end{table}

A previsibilidade do FSM é uma propriedade intencional: ela garante que o
\textit{baseline} forneça um referencial estável e reprodutível para a
comparação com o agente treinado por PPO \cite{millington2019ai}.

\subsection{Dano do Boss --- BossAttackDamage}
\label{subsec:boss-attack-damage}

O componente \texttt{BossAttackDamage} é compartilhado pelas duas condições
experimentais (FSM e PPO) e aplica dano real ao jogador quando um ataque do
chefe é executado dentro do alcance. A cada evento de ataque, o componente
verifica a distância até o jogador; se esta for inferior a
\texttt{attackRange} (1,25 unidades) e o \textit{cooldown} de dano (0,8
segundos) tiver expirado, 10 pontos de dano são aplicados ao
\texttt{PlayerHealth}.

O \textit{cooldown} de dano é especialmente relevante na condição PPO: como
o agente treinado apresenta alta frequência de ações de ataque, a ausência
de \textit{cooldown} resultaria em dano instantâneo e inviabilizaria a
avaliação. O valor de 0,8 segundos é idêntico ao \texttt{attackCooldown}
do FSM, preservando a equivalência de condições.

% -----------------------------------------------------------------------------
\section{Núcleo de Proveniência}
\label{sec:nucleo-proveniencia}

O sistema de proveniência é composto por três componentes que, em conjunto,
registram, armazenam e exportam as cadeias causais de eventos de cada sessão
de jogo.

\subsection{ProvenanceGraph --- Estrutura em Memória}
\label{subsec:provenance-graph}

O \texttt{ProvenanceGraph} mantém o grafo de proveniência da sessão corrente
em memória. Cada nó do grafo é um evento com os seguintes atributos:

\begin{itemize}
    \item \texttt{eventId}: identificador único do evento;
    \item \texttt{timestamp}: marcação temporal em segundos desde o início da
          sessão;
    \item \texttt{actorId}: identificador do ator que originou o evento
          (jogador ou chefe);
    \item \texttt{actionType}: tipo da ação (e.g., \texttt{PlayerAttack},
          \texttt{BossAttack}, \texttt{BossDamageTaken});
    \item \texttt{position}: posição 2D do ator no momento do evento;
    \item \texttt{value}: valor numérico associado (e.g., quantidade de dano);
    \item \texttt{parentEventId}: referência ao evento causador.
\end{itemize}

O atributo \texttt{parentEventId} é o elemento que distingue este sistema de
um log sequencial convencional. Ao apontar para o evento que causou o evento
corrente, ele estabelece uma relação causal explícita, transformando a
sequência temporal de eventos em um grafo dirigido acíclico (DAG) com
semântica de causa e efeito \cite{provenance_base}.

\subsection{ProvenanceLogger --- Captura de Eventos}
\label{subsec:provenance-logger}

O \texttt{ProvenanceLogger} atua como ponto de integração entre o
\textit{gameplay} e o sistema de proveniência. Ele inscreve-se nos eventos
emitidos pelos componentes de jogo e, para cada evento recebido, cria um nó
correspondente no \texttt{ProvenanceGraph} com o \texttt{parentEventId}
apontando para o evento mais recente na cadeia.

Os eventos capturados pelo \texttt{ProvenanceLogger} incluem:

\begin{itemize}
    \item \textbf{Ações do jogador}: \texttt{PlayerJump},
          \texttt{PlayerAttack}, \texttt{PlayerDash} --- emitidos pelo
          \texttt{PlayerController};
    \item \textbf{Dano no jogador}: \texttt{PlayerDamageTaken} --- emitido
          pelo \texttt{PlayerHealth} ao receber dano;
    \item \textbf{Ações do chefe}: \texttt{BossAttack} --- emitido pelo
          \texttt{BossFsmController} (ou equivalente no PPO);
    \item \textbf{Dano no chefe}: \texttt{BossDamageTaken} --- emitido pelo
          \texttt{BossHealth};
    \item \textbf{Morte do chefe}: \texttt{BossDeath} --- emitido pelo
          \texttt{BossHealth} quando a vida atinge zero;
    \item \textbf{Controle de sessão}: \texttt{SessionStart} e
          \texttt{SessionEnd} --- emitidos pelo \texttt{ArenaManager}.
\end{itemize}

Um aspecto relevante de projeto é que o \texttt{ProvenanceLogger} não modifica
nem interfere no fluxo de \textit{gameplay}. Ele é um observador puro: escuta
eventos, registra nós no grafo e não emite ações de volta ao jogo. Essa
separação garante que a presença ou ausência do sistema de proveniência não
altere o comportamento da simulação, propriedade essencial para a validade da
comparação entre as condições FSM e PPO.

A \autoref{fig:cadeia-causal} ilustra um exemplo de cadeia causal típica
registrada pelo sistema.

 \begin{figure}[htb]
     \centering
     \includegraphics[width=0.8\textwidth]{imagens/cadeia-causal}
     \caption{Exemplo de cadeia causal registrada pelo sistema de
              proveniência em uma sessão de combate.}
     \label{fig:cadeia-causal}
     \fonte{Elaborado pelo autor.}
 \end{figure}

A cadeia causal típica de uma sessão vitoriosa segue o encadeamento:

\begin{quote}
    \texttt{SessionStart} $\rightarrow$ \texttt{BossAttack} $\rightarrow$
    \texttt{PlayerDamageTaken} $\rightarrow$ \texttt{PlayerAttack}
    $\rightarrow$ \texttt{PlayerDamageDealt} $\rightarrow$
    \texttt{BossDamageTaken} $\rightarrow$ \texttt{BossDeath} $\rightarrow$
    \texttt{SessionEnd}
\end{quote}

Nessa cadeia, a relação entre \texttt{PlayerAttack} e
\texttt{BossDamageTaken} não é meramente temporal: o atributo
\texttt{parentEventId} do nó \texttt{BossDamageTaken} aponta explicitamente
para o nó \texttt{PlayerAttack}, indicando causalidade direta.

\subsection{ProvenanceExporter --- Serialização em JSON}
\label{subsec:provenance-exporter}

O \texttt{ProvenanceExporter} é responsável por serializar o grafo de
proveniência em formato JSON ao término de cada sessão. A exportação é
disparada pelo evento \texttt{SessionEnd} emitido pelo
\texttt{ArenaManager}, e o arquivo resultante é gravado no diretório
\texttt{TreinamentoML/provenance\_data/} com nome baseado no identificador
da sessão.

Cada arquivo JSON contém a lista completa de nós do grafo, incluindo todos os
atributos descritos na \autoref{subsec:provenance-graph}, além de metadados da
sessão como resultado (\texttt{victory} ou \texttt{defeat}), duração e tipo
de chefe. Esse formato permite que o script Python de extração de sequências
opere diretamente sobre os arquivos exportados, sem necessidade de
transformações intermediárias.

% -----------------------------------------------------------------------------
\section{Integração com Aprendizado por Reforço}
\label{sec:integracao-rl}

A integração entre o sistema de proveniência e o aprendizado por reforço é
realizada por dois componentes complementares: o
\texttt{ProvenanceRewardShaper}, que converte as sequências vitoriosas em
sinais de recompensa, e o \texttt{BossAgent}, que implementa o agente PPO
com observações vetoriais e ações discretas.

\subsection{ProvenanceRewardShaper --- Recompensa por Proveniência}
\label{subsec:reward-shaper-impl}

O \texttt{ProvenanceRewardShaper} é o componente que materializa a conexão
entre proveniência e aprendizado por reforço. Seu funcionamento segue três
etapas:

\begin{enumerate}
    \item \textbf{Carregamento}: na inicialização, o componente lê o arquivo
          \texttt{winning\_sequences.json} --- gerado pelo script Python de
          extração --- e armazena as sequências vitoriosas em memória como
          listas de \textit{strings};
    \item \textbf{Registro}: a cada ação executada pelo agente, a
          \textit{string} correspondente à ação (e.g.,
          \texttt{"PlayerAttack"}) é inserida em um \textit{buffer} circular
          de tamanho fixo (igual ao tamanho das sequências, i.e., 3);
    \item \textbf{Comparação}: após cada inserção, o conteúdo do
          \textit{buffer} é comparado com todas as sequências vitoriosas
          carregadas. Quando há correspondência exata, o componente retorna
          uma recompensa configurável com valor padrão de $+2{,}0$.
\end{enumerate}

O \texttt{BossAgent} consulta o \texttt{ProvenanceRewardShaper} a cada passo
de decisão e, caso haja recompensa disponível, a adiciona ao sinal de
recompensa do episódio via \texttt{AddReward}. Essa recompensa auxiliar atua
em conjunto com a recompensa extrínseca do ambiente, sem substituí-la.

A decisão de isolar a lógica de \textit{Reward Shaping} em um componente
separado foi motivada por quatro razões:

\begin{enumerate}
    \item Permite testar o mecanismo de correspondência de sequências
          isoladamente, sem depender do ciclo de treinamento do ML-Agents;
    \item Mantém o código do \texttt{BossAgent} focado em observações e
          ações, sem misturar lógica de proveniência;
    \item Facilita a substituição ou extensão do esquema de \textit{Reward
          Shaping} em trabalhos futuros;
    \item Preserva a rastreabilidade da contribuição acadêmica: toda a lógica
          que conecta proveniência e RL está localizada em um único componente
          auditável.
\end{enumerate}

\subsection{Vocabulário Compartilhado de Ações}
\label{subsec:vocabulario-acoes}

Um aspecto de implementação que merece destaque é a correspondência de
vocabulário entre as ações do agente PPO e as ações humanas registradas pela
proveniência. Quando o \texttt{BossAgent} executa uma ação de pulo, ataque ou
\textit{dash}, ele registra no \texttt{ProvenanceRewardShaper} a mesma
\textit{string} utilizada pelo \texttt{ProvenanceLogger} para a ação
equivalente do jogador humano:

\begin{itemize}
    \item Ação de pulo do agente $\rightarrow$ \texttt{"PlayerJump"}
    \item Ação de ataque do agente $\rightarrow$ \texttt{"PlayerAttack"}
    \item Ação de \textit{dash} do agente $\rightarrow$ \texttt{"PlayerDash"}
\end{itemize}

Essa correspondência garante que a comparação entre o \textit{buffer} do
agente e as sequências vitoriosas extraídas dos grafos de proveniência seja
direta, sem necessidade de mapeamentos intermediários. O agente PPO, ao
executar a mesma combinação de ações que um jogador humano utilizou para
causar dano no chefe em uma sessão vitoriosa, recebe a recompensa auxiliar.

\subsection{BossAgent --- Agente PPO}
\label{subsec:boss-agent-impl}

O \texttt{BossAgent} herda da classe \texttt{Agent} do Unity ML-Agents
\cite{juliani2020unity} e implementa os três métodos fundamentais do ciclo de
treinamento: \texttt{CollectObservations}, \texttt{OnActionReceived} e
\texttt{OnEpisodeBegin}.

\subsubsection{Espaço de Observação}

O agente coleta 10 observações vetoriais a cada passo de decisão, conforme
detalhado na \autoref{tab:observacoes}.

\begin{table}[htb]
    \centering
    \caption{Observações vetoriais do \texttt{BossAgent}.}
    \label{tab:observacoes}
    \begin{tabular}{r l p{6.5cm}}
        \toprule
        \textbf{\#} & \textbf{Observação}     & \textbf{Descrição} \\
        \midrule
        1  & \texttt{bossPosition.x}           & Posição horizontal do chefe. \\
        2  & \texttt{bossPosition.y}           & Posição vertical do chefe. \\
        3  & \texttt{bossVelocity.x}           & Velocidade horizontal do chefe. \\
        4  & \texttt{bossVelocity.y}           & Velocidade vertical do chefe. \\
        5  & \texttt{relativePlayerPos.x}      & Distância relativa horizontal
                                                 até o jogador. \\
        6  & \texttt{relativePlayerPos.y}      & Distância relativa vertical
                                                 até o jogador. \\
        7  & \texttt{normalizedBossHealth}      & Vida normalizada do chefe
                                                 (intervalo $[0, 1]$). \\
        8  & \texttt{1.0} (\textit{placeholder}) & Vida do jogador.
                                                 \textit{Placeholder} fixo
                                                 utilizado durante o treino. \\
        9  & \texttt{IsGrounded()}              & Indicador binário: 1 se o
                                                 chefe está no chão, 0 caso
                                                 contrário. \\
        10 & \texttt{0.0} (\textit{placeholder}) & \textit{Cooldown} de
                                                 \textit{dash}.
                                                 \textit{Placeholder} fixo. \\
        \bottomrule
    \end{tabular}
    \fonte{Elaborado pelo autor.}
\end{table}

As observações 8 e 10 são \textit{placeholders} fixos que foram preservados
na arquitetura final do modelo. Embora não forneçam informação variável ao
agente, sua presença no vetor de observação era necessária para manter a
compatibilidade dimensional com o modelo ONNX exportado. Essa decisão é uma
consequência direta do congelamento do modelo antes da adição do
\texttt{PlayerHealth} real, conforme documentado nas decisões de escopo
(\autoref{sec:decisoes-escopo}).

\subsubsection{Espaço de Ação}

O agente opera com um espaço de ação discreto contendo seis ações mutuamente
exclusivas, apresentadas na \autoref{tab:acoes-ppo}.

\begin{table}[htb]
    \centering
    \caption{Ações discretas do \texttt{BossAgent}.}
    \label{tab:acoes-ppo}
    \begin{tabular}{r l p{7cm}}
        \toprule
        \textbf{Índice} & \textbf{Ação}          & \textbf{Efeito} \\
        \midrule
        0 & \textit{Idle}          & Nenhuma ação. O chefe permanece parado. \\
        1 & \textit{Move Left}     & Desloca o chefe para a esquerda. \\
        2 & \textit{Move Right}    & Desloca o chefe para a direita. \\
        3 & \textit{Jump}          & Aplica impulso vertical ao chefe. Registra
                                     \texttt{"PlayerJump"} no \textit{Reward
                                     Shaper}. \\
        4 & \textit{Attack Melee}  & Executa ataque corpo a corpo. Registra
                                     \texttt{"PlayerAttack"} no \textit{Reward
                                     Shaper}. \\
        5 & \textit{Dash}          & Executa \textit{dash} horizontal. Registra
                                     \texttt{"PlayerDash"} no \textit{Reward
                                     Shaper}. \\
        \bottomrule
    \end{tabular}
    \fonte{Elaborado pelo autor.}
\end{table}

As ações 3, 4 e 5 correspondem ao vocabulário de ações que o
\texttt{ProvenanceRewardShaper} monitora. As ações 0, 1 e 2 não participam
do mecanismo de \textit{Reward Shaping}, pois não fazem parte do vocabulário
de ações registrado nos grafos de proveniência.

\subsubsection{Integração do Modelo ONNX}

O \textit{run} final de treinamento (\texttt{boss\_v1\_500k}) produziu o
modelo \texttt{BossAgent.onnx} após 500.000 passos de treinamento. As
métricas finais registradas pelo TensorBoard foram: recompensa média
acumulada de $1766{,}4$ com desvio-padrão de $12{,}8$, obtidas em
aproximadamente 2.008 segundos de treinamento.

O modelo foi integrado ao projeto Unity no diretório
\texttt{Assets/Models/BossAgent.onnx} e associado ao componente
\texttt{BossAgent} no \textit{prefab} \texttt{Boss\_PPO} em modo
\texttt{InferenceOnly}. Nesse modo, o agente utiliza o modelo ONNX para
tomar decisões a cada passo, sem conexão com o treinador Python. Essa
configuração garante que o comportamento observado durante a avaliação reflita
exclusivamente a política aprendida durante o treinamento, sem qualquer ajuste
\textit{online} \cite{onnx2019}.

% -----------------------------------------------------------------------------
\section{Sistema de Métricas e Gerenciamento de Sessão}
\label{sec:sistema-metricas}

O sistema de coleta de métricas opera como uma camada de instrumentação
acoplada aos eventos já existentes no \textit{gameplay}, sem interferir na
simulação.

\subsection{ArenaManager --- Ciclo de Vida da Sessão}
\label{subsec:arena-manager}

O \texttt{ArenaManager} coordena o ciclo de vida de cada sessão de combate.
Suas responsabilidades incluem:

\begin{itemize}
    \item Emitir o evento \texttt{SessionStart} ao iniciar a cena;
    \item Detectar vitória (morte do chefe via evento do \texttt{BossHealth})
          ou derrota (morte do jogador via evento do \texttt{PlayerHealth});
    \item Emitir o evento \texttt{SessionEnd} com o resultado
          (\texttt{victory} ou \texttt{defeat});
    \item Acionar a exportação de métricas CSV e, quando habilitada, a
          exportação do grafo de proveniência em JSON.
\end{itemize}

\subsection{GameMetrics e GameMetricsExporter --- Coleta e Exportação CSV}
\label{subsec:game-metrics}

O \texttt{GameMetrics} inscreve-se nos eventos de \textit{gameplay} e
acumula contadores e métricas ao longo da sessão: duração, dano recebido pelo
chefe, dano recebido pelo jogador, contagem de ações do jogador (pulos,
ataques, \textit{dashes}) e contagem de ações do chefe. Para sessões contra o
chefe PPO, são registradas também as contagens individuais de cada uma das
seis ações discretas do \texttt{BossAgent}.

Ao término da sessão, o \texttt{GameMetricsExporter} serializa todas as
métricas em uma linha CSV e a anexa ao arquivo
\texttt{boss\_evaluation\_metrics.csv}. O exportador implementa três
cuidados de engenharia:

\begin{enumerate}
    \item \textbf{Cabeçalho estável}: o cabeçalho é escrito uma única vez, na
          criação do arquivo, garantindo que todas as linhas compartilhem a
          mesma estrutura de colunas;
    \item \textbf{Formato invariante}: valores numéricos são formatados com
          \texttt{CultureInfo.InvariantCulture}, assegurando que o separador
          decimal seja sempre o ponto, independentemente da configuração
          regional do sistema operacional;
    \item \textbf{Resolução de caminho}: o caminho do arquivo é resolvido
          relativamente à raiz do repositório, e o diretório de destino é
          criado automaticamente caso não exista.
\end{enumerate}

Essas precauções garantem que o CSV gerado seja diretamente consumível por
ferramentas de análise sem necessidade de pré-processamento.
